\documentclass[10pt,a4paper,titlepage]{jreport} % ドキュメントクラスを1つに統一


\usepackage[top=25truemm,bottom=25truemm,left=20truemm,right=20truemm]{geometry} % 余白設定
\usepackage[dvipdfmx]{graphicx} % 画像の挿入
\usepackage{graphicx}
\usepackage{listings} % コードリスト用
\usepackage{jlisting} % 日本語対応のリスト用
\usepackage{jverb} % 日本語のベタ打ち
\usepackage{booktabs}
\usepackage{pgfplots} % グラフ描画用パッケージ
\pgfplotsset{compat=1.18} % バージョン指定
\usepackage{float}
\parindent = 0pt % 段落の字下げを無効化

\usepackage{titlesec}
\usepackage{etoolbox}
\makeatletter
\patchcmd{\chapter}{\if@openleft\cleardoublepage\else\if@openright\cleardoublepage\else\clearpage\fi\fi}{}{}{}
\makeatother
\lstset{breaklines=true, postbreak=\mbox{$\hookrightarrow$}\space, keepspaces=true, escapeinside={\%*}{*)}}

\titleformat{\chapter}[hang] % 章のフォーマットを変更
  {\normalfont\huge\bfseries} % フォント設定
  {\thechapter} % 番号部分の表示形式
  {1em} % 番号とタイトルの間のスペース
  {} % タイトルの前に入れる内容（ここでは空）
\usepackage{xcolor}
\usepackage{amsmath}
\usepackage{amssymb} % ←これで \mathbb{R} が使える

\lstset{
  language=Matlab,              % 言語はMATLAB
  basicstyle=\ttfamily\small,   % フォントは等幅、サイズ小さめ
  keywordstyle=\color{blue},    % キーワード（function, endとか）を青
  commentstyle=\color{green!50!black}, % コメントを深緑
  stringstyle=\color{red!70!black},    % 文字列（'...'）を赤系
  numbers=left,                 % 行番号を左に表示
  numberstyle=\tiny\color{gray},% 行番号をグレーで小さく
  stepnumber=1,                 % 毎行行番号つける
  numbersep=10pt,               % コードとの間隔
  backgroundcolor=\color{gray!10}, % 背景うっすらグレー
  frame=single,                 % 枠をつける
  rulecolor=\color{black},      % 枠線は黒
  breaklines=true,              % 長い行を折り返す
  breakatwhitespace=false,      % スペースでだけ折り返すか（falseにして自然な折り返し）
  showspaces=false,             % スペースは特別に表示しない
  showstringspaces=false,       % 文字列中のスペースも普通に表示
  tabsize=2                     % タブ幅2
}


\title{数値計算：課題2} % タイトル
\author{
  学生番号：242C2016  氏名：奥村直 \\
  \\
  知的システム工学科システム制御コース
  } % 著者
\date{\today} % 日付

\begin{document}

\maketitle

\section*{課題1}

\subsection*{問1}
IEEE754 規格の倍精度浮動小数点数において、
1 と $1+\varepsilon$ を区別できる最小の $\varepsilon$ を求める。

IEEE754 倍精度浮動小数点数は符号部 1 bit指数部 11 bit
仮数部 52 bit から構成され全体で 64 bit で表現される。
正規化表現では仮数の最上位ビットは常に $1$ となるため
数は次の形で表される。
\[
x = (1.d_2 d_3 \cdots d_{52})_2 \times 2^e
\]

数 $1$ は
\[
1 = 1.000\cdots000 \times 2^0
\]
と表される。これより 1 より大きい最小の表現可能な数は
\[
1.000\cdots001 \times 2^0
\]
でありその差は仮数部の最下位ビットに対応する値となる。

仮数部は 52 bit であるため
\[
\varepsilon = 2^{-52}
\]
となる。数値的には
\[
\varepsilon \approx 2.22 \times 10^{-16}
\]
である。

\subsection*{問2}
例1-1(2) において桁落ちが生じない方法を示し
それによる結果を求める。

次の式を考える。
\[
\sqrt{x+1} - \sqrt{x}
\]
$x$ が十分大きい場合$\sqrt{x+1}$ と $\sqrt{x}$ は非常に近い値となるため
直接減算を行うと有効桁が失われ桁落ちが生じる。

これを防ぐため分母分子に
$\sqrt{x+1} + \sqrt{x}$ を掛けて有理化を行う。
\[
\sqrt{x+1} - \sqrt{x}
= \frac{(\sqrt{x+1}-\sqrt{x})(\sqrt{x+1}+\sqrt{x})}{\sqrt{x+1}+\sqrt{x}}
\]

分子を整理すると
\[
= \frac{(x+1)-x}{\sqrt{x+1}+\sqrt{x}}
= \frac{1}{\sqrt{x+1}+\sqrt{x}}
\]

この変形により引き算を含まない形で計算できるため
桁落ちを回避することができる。


\end{document}